\documentclass[14pt]{extarticle}

\usepackage[
    a4paper,
    top=40mm,
    bottom=40mm,
    left=30mm,
    right=30mm
]{geometry}

\usepackage[macedonian]{babel}
\usepackage{graphicx}
\usepackage{amsmath}
\usepackage{amsfonts}
\usepackage{booktabs}
\usepackage{hyperref}

\setlength{\parindent}{0pt}
\setlength{\parskip}{0.6em}

\newcommand{\UniversityName}
    {Универзитет „Св. Кирил и Методиј“}
\newcommand{\FacultyName}
    {Факултет за информатички науки и компјутерско инженерство}
\newcommand{\ProjectType}
    {дипломска работа}
\newcommand{\ProjectTitle}
    {Напредни архитектури на граф невронски мрежи за пребарување и расудување над графови при одговарање на прашања над графови на знаење}
\newcommand{\ProjectMentor}
    {Име и презиме на менторот}
\newcommand{\ProjectCandidate}
    {Име и презиме на кандидатот}
\newcommand{\ProjectCandidateIndex}
    {Број на индекс}
\newcommand{\ProjectDate}
    {Месец и година}

\date{}

\begin{document}

\begin{titlepage}
    \begin{center}
        {\Large \UniversityName}\par
        \vspace{0.5cm}
        {\Large \FacultyName}\par

        \vspace*{\fill}

        {\Large\textbf{\MakeUppercase{\ProjectTitle}}}\par
        \vspace{1cm}
        {\large\textsc{\ProjectType}}\par

        \vspace*{\fill}

        \begin{minipage}[t]{0.49\textwidth}
            \large\textbf{Ментор}\par
            \ProjectMentor
        \end{minipage}
        \hfill
        \begin{minipage}[t]{0.45\textwidth}
            \raggedleft
            \large\textbf{Кандидат}\par
            \ProjectCandidate\par
            \ProjectCandidateIndex
        \end{minipage}

        \vspace{1cm}
        \rule{\textwidth}{0.5mm}
        {\large \ProjectDate}
    \end{center}
\end{titlepage}

\tableofcontents
\newpage

\listoftables
\listoffigures
\newpage

\section{Вовед}

\textbf{Мотивација.}

\textbf{Проблем и цел на трудот.}

\textbf{Истражувачки прашања.}

\textbf{Придонеси на трудот.}

\textbf{Организација на трудот.}

\section{Теоретска основа и сродни истражувања}

\textbf{Графови на знаење и одговарање на прашања.}

\textbf{Генерирање збогатено со пребарување.}

\textbf{Граф невронски мрежи.}

\textbf{Сродни архитектури.}

\textbf{Конструкција на доказен подграф.}

\textbf{Позиционирање на трудот.}

\section{Методологија и архитектура на системот}

\textbf{Формулација на задачата.}

\textbf{Преглед на целосниот тек.}

\textbf{Подготовка на WebQSP.}

\textbf{Заедничка рамка за пребарување на одговори.}

\textbf{GraphSAGE.}

\textbf{Напреден GraphSAGE.}

\textbf{R-GCN.}

\textbf{HGT.}

\textbf{ReaRev.}

\textbf{NBFNet.}

\textbf{Унија на најкратки патеки.}

\textbf{Prize-Collecting Steiner Tree.}

\textbf{Генерирање на конечниот одговор.}

\textbf{Репродуктивност.}

\section{Експериментална поставеност}

\textbf{Податочно множество.}

\textbf{Заедничка конфигурација.}

\textbf{Метрики за пребарување.}

\textbf{Метрики за доказниот контекст.}

\textbf{Метрики за конечниот одговор.}

\textbf{Експеримент I -- споредба на архитектурите.}

\textbf{Експеримент II -- споредба на доказните подграфови.}

\textbf{Експеримент III -- споредба на јазичните модели.}

\section{Резултати и дискусија}

\textbf{Споредба на архитектурите за пребарување.}

\textbf{Влијание на конструкцијата на доказниот подграф.}

\textbf{Влијание на јазичниот модел.}

\textbf{Загуба на информации низ системот.}

\textbf{Ограничувања и закани по валидноста.}

\section{Заклучок}

\section{Идни подобрувања}

\bibliographystyle{plain}
\bibliography{references}

\end{document}
